%--------------------------------------------------------------------%
% BAB I MASALAH MITRA DAN TUJUAN PROYEK
%--------------------------------------------------------------------%
\chapter{MASALAH MITRA DAN TUJUAN PROYEK} \label{bab:1}
\petunjuk{Bab ini menjelaskan masalah nyata mitra dan tujuan yang disepakati, bukan latar belakang penelitian. Subbab 1.1 sampai 1.4 boleh beririsan dengan laporan anggota tim lain. \tahap{\wajib{WAJIB (draf)}}{\wajib{WAJIB (final)}}}

\section{Profil Mitra} \label{sec:1.1}
\petunjuk{Ringkasan profil mitra dari P.1: bidang usaha atau layanan dan proses yang menjadi konteks proyek.}
\tulisuraian

\section{Masalah Mitra} \label{sec:1.2}
\petunjuk{Selisih antara kondisi saat ini dan kondisi yang diharapkan mitra untuk setiap aspek, dengan bukti (data, wawancara, atau dokumen mitra).}
\tulisuraian

\section{Tujuan dan Kriteria Penerimaan} \label{sec:1.3}
\petunjuk{Untuk setiap tujuan, uraikan kriteria penerimaan dan cara verifikasinya. Tujuan dan kriteria disepakati dengan mitra dan harus sama dengan proposal; perubahan didokumentasikan dan disetujui pembimbing. \tahap{\wajib{WAJIB}}{\wajib{WAJIB (final)}}}
\tulisuraian

\section{Ruang Lingkup dan Batasan} \label{sec:1.4}
\petunjuk{Apa yang dikerjakan dan yang tidak dikerjakan dalam proyek.}
\tulisuraian

\section{Bagian Individu dan Posisinya dalam Proyek} \label{sec:1.5}
\petunjuk{Nyatakan judul laporan ini, bagian proyek yang menjadi tanggung jawab Anda sesuai Surat Pembagian Kerja Tim, keterkaitannya dengan bagian anggota lain (tandai bagian Anda pada diagram arsitektur umum), serta tujuan dan kriteria keberhasilan bagian Anda. Bagian ini wajib ditulis sendiri dan berbeda dari laporan anggota tim lain (Pedoman TA 5.1.3). \tahap{\wajib{WAJIB}}{\wajib{WAJIB (final)}}}
\tulisuraian
